\NSPage{049}%
\begin{EESegment}{EE-TGT-P049-001}
equation across \NSi{NS-F-P049-001}{\xi=0}. The original equations keep the first four coordinates even and the last two
odd, and uniqueness means that these are exactly the parities of the solution. For an even smooth function of
\NSi{NS-F-P049-002}{\xi}, Taylor's formula shows at every finite order that the function is smooth in
\NSi{NS-F-P049-003}{X=\xi^2} on the half-interval. This also gives the regularity needed to form the next source
\NSi{NS-F-P049-004}{\Omega_{n}/X}.
\end{EESegment}
\end{proof}

\begin{EESegment}{EE-TGT-P049-002}
\subsection{Making the total radial residual integrals zero}\label{sec:5.2}
The inner profiles \NSi{NS-F-P049-005}{(\phi_n,U_n,V_n,\Pi_n)} from
Lemma~\ref{lem:5.1} are not yet ready to be used as coefficients at the next order. First, extend
\NSi{NS-F-P049-006}{E_n,U_n} to every \NSi{NS-F-P049-007}{X\geq 0}. Then recompute
\NSi{NS-F-P049-008}{V_n,\Pi_n} from Equations~\eqref{eq:5.2} and~\eqref{eq:5.5}, and control the residual that the
extension creates. Cutting off only \NSi{NS-F-P049-009}{E_n,U_n} in the radial direction does not make the fields found by
radial integration compactly supported. In particular, the pressure can be left with a nonzero constant outside the
support, and the stress \NSi{NS-F-P049-010}{\mathcal T_n} defined below in \eqref{eq:5.9} can be left with a nonzero tail
there. To remove these terms, impose five moment conditions, which here means setting the five integrals in
Equations~\eqref{eq:5.10} and~\eqref{eq:5.11} equal to zero. Lemma~\ref{lem:5.2} proves that these conditions do remove
the terms.
\end{EESegment}

\begin{EESegment}{EE-TGT-P049-003}
Here is how the residual coefficients are collected before the moment conditions are imposed. Use
\NSi{NS-F-P049-011}{R=\sqrt{2X}}. To find the part of order \NSi{NS-F-P049-012}{n}, substitute \eqref{eq:5.1} into the
differential polynomial \NSi{NS-F-P049-013}{\mathscr R}, take the physical derivatives, divide out the common tangential
power \NSi{NS-F-P049-014}{q^{-A-1}}, and keep the coefficient of \NSi{NS-F-P049-015}{q^{2nh}}. This calculation is
finite. Products use the pairs \NSi{NS-F-P049-016}{i+j=n}, while axial viscosity uses
order \NSi{NS-F-P049-017}{n-1}. More explicitly,
\NSFormula{NS-F-P049-018}{display}{none}
\[
\begin{aligned}
\mathfrak r_{\theta,n}
  &=\frac{R}{\mathsf C}\left[
      T_{b,n}\phi_n
      +\sum_{i+j=n}\left\{V_i\left(\phi_j'+\frac{\phi_j}{X}\right)+U_iZ_{b,j}\phi_j\right\}
      -2\left(X\phi_n''+2\phi_n'\right)
      -Z_{b,n-1}^{[2]}\phi_{n-1}
    \right],\\
\mathfrak r_{z,n}
  &=T_{c,n}U_n
      +\sum_{i+j=n}\left\{V_iU_j'+U_iZ_{c,j}U_j\right\}
      +Z_{-2A,n}\Pi_n
      -2\left(XU_n''+U_n'\right)
      -Z_{c,n-1}^{[2]}U_{n-1}.
\end{aligned}
\]
\end{EESegment}

\begin{EESegment}{EE-TGT-P049-004}
These expressions are zero wherever the inner equations have been solved. Define the two needed stress components by
integrating the negatives of the residuals with the indicated radial weights.
\NSFormula{NS-F-P049-019}{display}{5.9}
\begin{equation}\label{eq:5.9}
\mathcal T_{n,\theta}(R)=-R^{-2}\int_0^R\varrho^2\mathfrak r_{\theta,n}(\varrho)\,d\varrho,
\qquad
\mathcal T_{n,z}(R)=-R^{-1}\int_0^R\varrho\mathfrak r_{z,n}(\varrho)\,d\varrho.
\tag{5.9}
\end{equation}
\end{EESegment}

\begin{EESegment}{EE-TGT-P049-005}
Every profile inside these integrals uses the same value of \NSi{NS-F-P049-020}{\eta}. The two physical stress components
are \NSi{NS-F-P049-021}{q^{-A-1/2+\lambda_n}\mathcal T_n}. For each forward integral to become zero after
\NSi{NS-F-P049-022}{R^2\mathfrak r_{\theta,n}} or \NSi{NS-F-P049-023}{R\mathfrak r_{z,n}} leaves the support of its
source, the corresponding total integral must be zero. The five conditions below force both total integrals to be zero
through the conservative momentum equations. They also remove the radial-velocity and pressure terms left by integration.
Also write \NSi{NS-F-P049-024}{F_n=\int_0^XU_n(x,\eta)\,dx=X\mathcal A_X(U_n)} for the Stokes streamfunction
coefficient. This indexed notation is different from the leading azimuthal profile
\NSi{NS-F-P049-025}{F=E/\sqrt{2X}}.
\end{EESegment}

\begin{EESegment}{EE-TGT-P049-006}
For an extension at order \NSi{NS-F-P049-026}{n} that has not yet been corrected, call the five total moments
\NSFormula{NS-F-P049-027}{display}{5.10}
\begin{equation}\label{eq:5.10}
m_{n,1}=\int_0^\infty RU_n\,dR,
\qquad
m_{n,2}=\int_0^\infty R^2E_n\,dR,
\qquad
m_{n,3}=\int_0^\infty \partial_R\Pi_n\,dR.
\tag{5.10}
\end{equation}
\NSFormula{NS-F-P049-028}{display}{5.11}
\begin{equation}\label{eq:5.11}
\begin{aligned}
m_{n,4}&=\int_0^\infty R^2\sum_{i+j=n}U_iE_j\,dR,\\
m_{n,5}&=\int_0^\infty\left(R\sum_{i+j=n}U_iU_j-\frac12R^2\partial_R\Pi_n\right)\,dR.
\end{aligned}
\tag{5.11}
\end{equation}
\end{EESegment}

\NSPage{050}%
\begin{EESegment}{EE-TGT-P050-001}
These five moments depend on \NSi{NS-F-P050-001}{\eta}, with every lower-order factor kept fixed. Write
\NSFormula{NS-F-P050-002}{display}{none}
\[
\mathbf m_n=(m_{n,1},\ldots,m_{n,5}),
\qquad
\mathbf P_n=\left(\phi_n,U_n,\frac{V_n}{X},\Pi_n,\frac{F_n}{X}\right).
\]
Every radial support used below is uniform for \NSi{NS-F-P050-003}{|\eta|\leq 1}. The stress estimates use the flat
weight \NSi{NS-F-P050-004}{\zeta} from Item~\textup{(iv)} of Theorem~\ref{thm:4.6}, together with
\NSFormula{NS-F-P050-005}{display}{none}
\[
\delta(X)=\min\{1,\log(X/X_a),\log(X_b/X)\}
\qquad (X_a<X<X_b).
\]
\end{EESegment}

\begin{lemma}\label{lem:5.2}
\begin{EESegment}{EE-TGT-P050-002}
There are radii \NSi{NS-F-P050-006}{X_a<X_-<X_+<X_b}, chosen independently of
\NSi{NS-F-P050-007}{n}, for the following induction. Fix \NSi{NS-F-P050-008}{n\geq 1}. For every
\NSi{NS-F-P050-009}{1\leq j<n}, suppose that \NSi{NS-F-P050-010}{\mathbf P_j} is smooth, satisfies
Equations~\eqref{eq:5.2} and~\eqref{eq:5.5} everywhere, and obeys \eqref{eq:5.12} with
\NSi{NS-F-P050-011}{n=j}. Suppose as well that these profiles agree with their inner solutions on
\NSi{NS-F-P050-012}{[0,X_-]} and keep the parameter regularity from Lemma~\ref{lem:5.1} throughout
\NSi{NS-F-P050-013}{[0,a^2]}. When \NSi{NS-F-P050-014}{n=1}, there are no lower positive-order hypotheses.

Then the inner order-\NSi{NS-F-P050-015}{n} solution from Lemma~\ref{lem:5.1} extends smoothly to every
\NSi{NS-F-P050-016}{X\geq0} and \NSi{NS-F-P050-017}{|\eta|\leq1}. The extension agrees with the inner solution on
\NSi{NS-F-P050-018}{[0,X_-]}, keeps its parameter regularity throughout \NSi{NS-F-P050-019}{[0,a^2]}, satisfies
Equations~\eqref{eq:5.2} and~\eqref{eq:5.5} everywhere, and obeys
\NSFormula{NS-F-P050-020}{display}{5.12}
\begin{equation}\label{eq:5.12}
\mathbf m_n(\eta)=0\quad (|\eta|\leq1),
\qquad
\supp_X\mathbf P_n\subset[0,X_+].
\tag{5.12}
\end{equation}
Both \NSi{NS-F-P050-021}{V_n/X} and \NSi{NS-F-P050-022}{F_n/X} extend smoothly to
\NSi{NS-F-P050-023}{X=0}. The stress defined in \eqref{eq:5.9} satisfies
\NSFormula{NS-F-P050-024}{display}{none}
\[
\supp_X\mathcal T_1\subset[X_-,X_b],
\qquad
\supp_X\mathcal T_n\subset[X_-,X_+]\quad(n\geq2),
\]
and every fixed profile derivative \NSi{NS-F-P050-025}{\partial^I} satisfies
\NSFormula{NS-F-P050-026}{display}{5.13}
\begin{equation}\label{eq:5.13}
|\partial^I\mathcal T_n(X,\eta)|
\leq C_{n,I}\zeta(X)\delta(X)^{-N_{n,I}}
\qquad (X_a<X<X_b,\ |\eta|\leq1).
\tag{5.13}
\end{equation}
The constants \NSi{NS-F-P050-027}{C_{n,I},N_{n,I}} may depend on the order and the derivative. When
\NSi{NS-F-P050-028}{n\geq2}, it is possible to take \NSi{NS-F-P050-029}{N_{n,I}=0}.
\end{EESegment}
\end{lemma}

\begin{proof}
\begin{EESegment}{EE-TGT-P050-003}
\textbf{Step 1. Extend the current profiles without changing the inner solution.}
Keep the inner solution unchanged on \NSi{NS-F-P050-030}{[0,X_-]}, then choose the rest of the extension so that
\NSi{NS-F-P050-031}{\mathbf m_n(\eta)=0} for every \NSi{NS-F-P050-032}{|\eta|\leq1}. Setting
\NSi{NS-F-P050-033}{m_{n,1}=m_{n,3}=0} removes the streamfunction and pressure terms outside the support. Setting
\NSi{NS-F-P050-034}{m_{n,1}=m_{n,2}=0} cancels the integrated time and viscosity terms. Setting
\NSi{NS-F-P050-035}{m_{n,4}=m_{n,5}=0} cancels the integrated momentum fluxes, including the pressure term.
Steps 3 and 4 show why each of these statements follows.
\end{EESegment}

\begin{EESegment}{EE-TGT-P050-004}
The inner profiles
\NSi{NS-F-P050-036}{(\phi_n^{\mathrm{in}},U_n^{\mathrm{in}},V_n^{\mathrm{in}},\Pi_n^{\mathrm{in}})} are defined for
\NSi{NS-F-P050-037}{0\leq X\leq a^2} by Lemma~\ref{lem:5.1}. Let
\NSi{NS-F-P050-038}{\mathcal I_{\mathrm{pos}}} be the reserved positive-order patch from Item~\textup{(vi)} of
Theorem~\ref{thm:4.6}. Choose \NSi{NS-F-P050-039}{X_-} inside the theorem's inner collar, where the stress cone
inequalities have a positive margin that is uniform over the collar. Then fix, independently of
\NSi{NS-F-P050-040}{n},
\NSFormula{NS-F-P050-041}{display}{none}
\[
X_-<X_{\mathrm{keep}}<X_{\mathrm{cut}}<a^2<\inf\mathcal I_{\mathrm{pos}},
\qquad
\kappa\in C^\infty([0,\infty)),
\qquad
0\leq\kappa\leq1,
\]
\NSFormula{NS-F-P050-042}{display}{none}
\[
\kappa(X)=1\quad(0\leq X\leq X_{\mathrm{keep}}),
\qquad
\kappa(X)=0\quad(X\geq X_{\mathrm{cut}}).
\]
\end{EESegment}

\begin{EESegment}{EE-TGT-P050-005}
Start with the extensions
\NSFormula{NS-F-P050-043}{display}{none}
\[
\widetilde\phi_n=\kappa\phi_n^{\mathrm{in}},
\qquad
\widetilde U_n=\kappa U_n^{\mathrm{in}},
\qquad
\widetilde E_n=\frac{R}{\mathsf C}\widetilde\phi_n,
\qquad
R=\sqrt{2X},
\]
and extend the products by zero past \NSi{NS-F-P050-044}{X=a^2}.
\end{EESegment}

\begin{EESegment}{EE-TGT-P050-006}
In the \NSi{NS-F-P050-045}{R} coordinate, the reserved patch \NSi{NS-F-P050-046}{\mathcal I_{\mathrm{pos}}} becomes
\NSi{NS-F-P050-047}{\mathcal J_{\mathrm{pos}}=\{R>0:R^2/2\in\mathcal I_{\mathrm{pos}}\}}. Choose the following
nonnegative smooth bump functions. Their supports are mutually disjoint, appear in a fixed order, and stay the same at
every order.
\NSFormula{NS-F-P050-048}{display}{none}
\[
b_1^U,b_2^U,b_1^E,b_2^E,b_3^E\in C_c^\infty(\mathcal J_{\mathrm{pos}}),
\qquad
\int_0^\infty b_j^U\,dR=1\quad(j=1,2),
\qquad
\int_0^\infty b_j^E\,dR=1\quad(j=1,2,3),
\]
\end{EESegment}

\NSPage{051}%
\begin{EESegment}{EE-TGT-P051-001}
Take five unknown scalar functions \NSi{NS-F-P051-001}{\alpha_{n,j}(\eta),\beta_{n,j}(\eta)} and set
\NSFormula{NS-F-P051-002}{display}{5.14}
\begin{equation}\label{eq:5.14}
\begin{aligned}
U_n(X,\eta)&=\widetilde U_n(X,\eta)+\sum_{j=1}^{2}\alpha_{n,j}(\eta)b_j^U(R),\\
E_n(X,\eta)&=\widetilde E_n(X,\eta)+\sum_{j=1}^{3}\beta_{n,j}(\eta)b_j^E(R),
\qquad
\phi_n=\frac{\mathsf C}{R}E_n.
\end{aligned}
\tag{5.14}
\end{equation}
At \NSi{NS-F-P051-003}{R=0}, use the smooth extension of the quotient that defines
\NSi{NS-F-P051-004}{\phi_n}. The bump functions vanish near the axis, where
\NSi{NS-F-P051-005}{\phi_n=\widetilde\phi_n}.
\end{EESegment}

\begin{EESegment}{EE-TGT-P051-002}
For any choice of these five coefficient functions, define the remaining profiles by
\NSFormula{NS-F-P051-006}{display}{5.15}
\begin{equation}\label{eq:5.15}
\begin{aligned}
F_n(X,\eta)&=\int_0^XU_n(x,\eta)\,dx,\\
V_n(X,\eta)&=-\int_0^X(Z_{-A,n}U_n)(x,\eta)\,dx,\\
\Pi_n(X,\eta)&=\int_0^X\left[\mathsf C^{-2}\sum_{i+j=n}\phi_i\phi_j-\frac{\Omega_{n-1}}{2x}\right](x,\eta)\,dx.
\end{aligned}
\tag{5.15}
\end{equation}
Every lower-order profile and \NSi{NS-F-P051-007}{\Omega_{n-1}} stays fixed in these formulas. The quotient
\NSi{NS-F-P051-008}{\Omega_{n-1}/x} is smooth at the axis, as shown after \eqref{eq:5.6}. Equations~\eqref{eq:5.2}
and~\eqref{eq:5.5} hold everywhere by construction, because the three profiles in \eqref{eq:5.15} were defined by
integrating those equations from the axis.
\end{EESegment}

\begin{EESegment}{EE-TGT-P051-003}
On \NSi{NS-F-P051-009}{[0,X_-]}, the cutoff is one and every bump is zero. Integrating forward from the zero values at
the axis therefore gives
\NSFormula{NS-F-P051-010}{display}{none}
\[
(\phi_n,U_n,V_n,\Pi_n)
=(\phi_n^{\mathrm{in}},U_n^{\mathrm{in}},V_n^{\mathrm{in}},\Pi_n^{\mathrm{in}}).
\]
The radial cutoff does not depend on \NSi{NS-F-P051-011}{\eta}, and the bump functions vanish throughout
\NSi{NS-F-P051-012}{[0,a^2]}. The reconstructed profiles therefore keep the parameter regularity from
Lemma~\ref{lem:5.1} on that whole interval, which is what the next order needs.
\end{EESegment}

\begin{EESegment}{EE-TGT-P051-004}
\textbf{Step 2. Solve the five moment equations.}
On \NSi{NS-F-P051-013}{\mathcal I_{\mathrm{pos}}},
\NSFormula{NS-F-P051-014}{display}{none}
\[
U_0=0,
\qquad
E_0=e_*f(\eta)R^{-1-2\lambda},
\qquad
e_*>0,
\qquad
\lambda>0,
\]
where \NSi{NS-F-P051-015}{f} is smooth and stays a positive distance away from zero. Once the lower-order profiles have
been fixed, each of the five moments is an affine function of the five unknown coefficients. For the fifth moment, the
pressure equation gives
\NSFormula{NS-F-P051-016}{display}{none}
\[
m_{n,5}=\int_0^\infty\left(
\sum_{i+j=n}U_iU_j
-\frac12\sum_{i+j=n}E_iE_j
+\frac12\Omega_{n-1}
\right)\,dX.
\]
\end{EESegment}

\begin{EESegment}{EE-TGT-P051-005}
The term \NSi{NS-F-P051-017}{\Omega_{n-1}} depends only on lower orders, so none of the five coefficients changes it.
Using \NSi{NS-F-P051-018}{R^2=2X}, the pressure equation also gives
\NSFormula{NS-F-P051-019}{display}{none}
\[
\partial_R\Pi_n
=\frac{2E_0E_n}{R}
+\frac1R\left(\sum_{i=1}^{n-1}E_iE_{n-i}-\Omega_{n-1}\right).
\]
\end{EESegment}

\begin{EESegment}{EE-TGT-P051-006}
Define the constant moment matrices
\NSFormula{NS-F-P051-020}{display}{none}
\[
\begin{aligned}
(B_U)_{ij}&=\int_0^\infty R^{p_i}b_j^U(R)\,dR,
& (p_1,p_2)&=(1,1-2\lambda),\\
(B_E)_{ij}&=\int_0^\infty R^{s_i}b_j^E(R)\,dR,
& (s_1,s_2,s_3)&=(2,-2-2\lambda,-2\lambda).
\end{aligned}
\]
\end{EESegment}

\NSPage{052}%
\begin{EESegment}{EE-TGT-P052-001}
The matrix \NSi{NS-F-P052-001}{B_U} is \NSi{NS-F-P052-002}{2\times2}, and
\NSi{NS-F-P052-003}{B_E} is \NSi{NS-F-P052-004}{3\times3}. Let \NSi{NS-F-P052-005}{m_{n,j}^0} be the moments obtained
by setting all five coefficients in \eqref{eq:5.14} equal to zero and then rebuilding the profiles with \eqref{eq:5.15}.
Put
\NSFormula{NS-F-P052-006}{display}{none}
\[
\mathbf d_{U,n}
=\begin{pmatrix}
m_{n,1}^0\\[2pt]
m_{n,4}^0/(e_*f)
\end{pmatrix},
\qquad
\mathbf d_{E,n}
=\begin{pmatrix}
m_{n,2}^0\\[2pt]
m_{n,3}^0/(2e_*f)\\[2pt]
-m_{n,5}^0/(e_*f)
\end{pmatrix}.
\]
\end{EESegment}

\begin{EESegment}{EE-TGT-P052-002}
With \NSi{NS-F-P052-007}{\boldsymbol\alpha_n=(\alpha_{n,1},\alpha_{n,2})^{\mathsf T}} and
\NSi{NS-F-P052-008}{\boldsymbol\beta_n=(\beta_{n,1},\beta_{n,2},\beta_{n,3})^{\mathsf T}}, the five moment conditions
say exactly that
\NSFormula{NS-F-P052-009}{display}{none}
\[
\mathbf m_n=0
\quad\Longleftrightarrow\quad
B_U\boldsymbol\alpha_n=-\mathbf d_{U,n},
\qquad
B_E\boldsymbol\beta_n=-\mathbf d_{E,n}.
\]
The exponents in each block are different, and the bump supports are ordered, so Lemma~\ref{lem:A.1} applies. Both
matrices are therefore invertible, and the needed coefficients are
\NSFormula{NS-F-P052-010}{display}{5.16}
\begin{equation}\label{eq:5.16}
\boldsymbol\alpha_n=-B_U^{-1}\mathbf d_{U,n},
\qquad
\boldsymbol\beta_n=-B_E^{-1}\mathbf d_{E,n}.
\tag{5.16}
\end{equation}
\end{EESegment}

\begin{EESegment}{EE-TGT-P052-003}
The discrepancies depend smoothly on \NSi{NS-F-P052-011}{\eta}, and every fixed derivative of
\NSi{NS-F-P052-012}{1/f} is bounded on \NSi{NS-F-P052-013}{[-1,1]}. The coefficients are therefore smooth whenever
the discrepancies have finite size, however large that size is. This proves \NSi{NS-F-P052-014}{\mathbf m_n=0}. The
support statement in \eqref{eq:5.12} still has to be proved.
\end{EESegment}

\begin{EESegment}{EE-TGT-P052-004}
\textbf{Step 3. Make the velocity and pressure zero outside the correction region.}
Choose \NSi{NS-F-P052-015}{X_+} after the leading axial-velocity perturbation \NSi{NS-F-P052-016}{U_0} and every
positive-order patch, but before the terminal outer collar. The first moment then gives
\NSFormula{NS-F-P052-017}{display}{none}
\[
F_n(X,\eta)=\int_0^\infty U_n(x,\eta)\,dx=m_{n,1}(\eta)=0
\qquad (X\geq X_+).
\]
So \NSi{NS-F-P052-018}{\mathcal A_X(U_n)=V_n=0} there at every positive order. At order zero, the leading
axial-flux moment correction also gives \NSi{NS-F-P052-019}{U_0=V_0=0} there. Every term in every
\NSi{NS-F-P052-020}{\Omega_k} contains either a \NSi{NS-F-P052-021}{V} factor or a derivative of a
\NSi{NS-F-P052-022}{V} field, so each such term is zero on the same outside region.
\end{EESegment}

\begin{EESegment}{EE-TGT-P052-005}
At positive order, every product \NSi{NS-F-P052-023}{\phi_i\phi_j} with \NSi{NS-F-P052-024}{i+j=n} contains a
positive-order factor. It follows that \NSi{NS-F-P052-025}{\partial_X\Pi_n=0} after
\NSi{NS-F-P052-026}{X_+}. The third moment now gives
\NSFormula{NS-F-P052-027}{display}{none}
\[
\Pi_n(X,\eta)
=\Pi_n(0,\eta)+\int_0^\infty\partial_R\Pi_n(R,\eta)\,dR
=m_{n,3}(\eta)=0
\qquad (X\geq X_+).
\]
Thus \NSi{NS-F-P052-028}{\phi_n,U_n,V_n,\Pi_n,F_n} are all zero for
\NSi{NS-F-P052-029}{X\geq X_+}. They are smooth because the cutoff construction is smooth, the solve in
\eqref{eq:5.16} is finite, and the remaining profiles come from the integrals in \eqref{eq:5.15}. At the axis,
\NSFormula{NS-F-P052-030}{display}{none}
\[
\frac{F_n(X,\eta)}{X}=\int_0^1U_n(tX,\eta)\,dt
\]
is smooth, and \eqref{eq:5.2} gives the same result for \NSi{NS-F-P052-031}{V_n/X}. Since
\NSi{NS-F-P052-032}{[0,X_+]\times[-1,1]} is compact, the coefficient bounds follow.
\NSFormula{NS-F-P052-033}{display}{5.17}
\begin{equation}\label{eq:5.17}
\max_{k+\ell\leq m}\ \sup_{X\geq0,\,|\eta|\leq1}
|\partial_X^k\partial_\eta^\ell\mathbf P_n(X,\eta)|
\leq C_{n,m}<\infty
\qquad(n\geq1,\ m\geq0).
\tag{5.17}
\end{equation}
The support radius is the same at every order, although these constants can grow with
\NSi{NS-F-P052-034}{n} and \NSi{NS-F-P052-035}{m}.
\end{EESegment}

\begin{EESegment}{EE-TGT-P052-006}
On the reserved interval \NSi{NS-F-P052-036}{\mathcal I_{\mathrm{mean}}}, which will later hold the angular-mean
corrections, the higher-order velocity coefficients are exactly zero. The inner cutoff and the patch
\NSi{NS-F-P052-037}{\mathcal I_{\mathrm{pos}}} both have their support to the left of
\NSi{NS-F-P052-038}{\mathcal I_{\mathrm{mean}}}. Thus \NSi{NS-F-P052-039}{E_n=U_n=0} there. The first moment makes
\NSi{NS-F-P052-040}{F_n=0} as soon as the positive-order patch has been passed, and \eqref{eq:5.2} then gives
\NSi{NS-F-P052-041}{V_n=0}. In particular,
\NSFormula{NS-F-P052-042}{display}{5.18}
\begin{equation}\label{eq:5.18}
E_n=U_n=F_n=V_n=0
\quad\text{on the reserved mean patch, for every }n\geq1.
\tag{5.18}
\end{equation}
It is essential to put \NSi{NS-F-P052-043}{X_+} after the leading axial-velocity perturbation if the pressure is to have
compact support, because \NSi{NS-F-P052-044}{\Omega_0} can still be nonzero after the positive-order patches.
\end{EESegment}

\NSPage{053}%
\begin{EESegment}{EE-TGT-P053-001}
\textbf{Step 4. Make the total tangential residual integrals zero.}
The inner equations make both positive-order stress components zero on \NSi{NS-F-P053-001}{[0,X_-]}. To show that
they are also zero after the outer radius, use the following conservative forms of the tangential residual, valid for any
smooth axisymmetric divergence-free field.
\NSFormula{NS-F-P053-002}{display}{none}
\[
\begin{aligned}
r^2\mathscr R_\theta
&=\partial_t(r^2u_\theta)
  +\partial_r(r^2u_ru_\theta)
  +\partial_z(r^2u_zu_\theta)
  -\partial_r(r^2\partial_ru_\theta-ru_\theta)
  -\partial_{zz}(r^2u_\theta),\\
r\mathscr R_z
&=\partial_t(ru_z)
  +\partial_r(ru_ru_z)
  +\partial_z\bigl(r(u_z^2+p)\bigr)
  -\partial_r(r\partial_ru_z)
  -\partial_{zz}(ru_z).
\end{aligned}
\]
\end{EESegment}

\begin{EESegment}{EE-TGT-P053-002}
The parity at the axis and the compact supports at positive order make the radial boundary terms zero. At order
\NSi{NS-F-P053-003}{n}, the angular moments that remain before differentiation are
\NSFormula{NS-F-P053-004}{display}{none}
\[
\begin{aligned}
\int_0^\infty r^2u_{\theta,n}\,dr
&=q^{3/2-A+\lambda_n}\int_0^\infty R^2E_n\,dR,\\
\int_0^\infty r^2\sum_{i+j=n}u_{z,i}u_{\theta,j}\,dr
&=q^{3/2-2A+\lambda_n}\int_0^\infty R^2\sum_{i+j=n}U_iE_j\,dR.
\end{aligned}
\]
Both are zero because \NSi{NS-F-P053-005}{m_{n,2}=m_{n,4}=0}. The power in the product is the same for every
split \NSi{NS-F-P053-006}{i+j=n}, so the cancellation happens before any derivative in the physical variables
\NSi{NS-F-P053-007}{z,t} is taken. The lower angular moment used by axial viscosity is zero because
\NSi{NS-F-P053-008}{m_{n-1,2}=0} when \NSi{NS-F-P053-009}{n\geq2}.
\end{EESegment}

\begin{EESegment}{EE-TGT-P053-003}
When \NSi{NS-F-P053-010}{n=1}, the compensated leading profile supplies the following convergent identity instead.
\NSFormula{NS-F-P053-011}{display}{5.19}
\begin{equation}\label{eq:5.19}
\int_0^\infty r^2\bigl(u_{\theta,0}(r,z,t)-P(r)\bigr)\,dr=0,
\tag{5.19}
\end{equation}
Here \NSi{NS-F-P053-012}{P} is the fixed pure-power exterior from \eqref{eq:A.45}, so
\NSi{NS-F-P053-013}{\partial_z^2P(r)=0}. Differentiation under the integral is valid. Beyond the moving terminal collar,
\NSi{NS-F-P053-014}{u_{\theta,0}=K(r,t)} does not depend on \NSi{NS-F-P053-015}{z}, and the difference moment itself
converges. The order-one angular viscosity term therefore has zero total moment as well.
\end{EESegment}

\begin{EESegment}{EE-TGT-P053-004}
The axial moments are
\NSFormula{NS-F-P053-016}{display}{5.20}
\begin{equation}\label{eq:5.20}
\begin{aligned}
\int_0^\infty ru_{z,n}\,dr
&=q^{1-A+\lambda_n}\int_0^\infty RU_n\,dR,\\
\int_0^\infty r\left(\sum_{i+j=n}u_{z,i}u_{z,j}+p_n\right)\,dr
&=q^{1-2A+\lambda_n}\int_0^\infty R\left(\sum_{i+j=n}U_iU_j+\Pi_n\right)\,dR.
\end{aligned}
\tag{5.20}
\end{equation}
Integration by parts turns the pressure part into
\NSFormula{NS-F-P053-017}{display}{none}
\[
\int_0^\infty R\Pi_n\,dR=-\frac12\int_0^\infty R^2\partial_R\Pi_n\,dR.
\]
The identities \NSi{NS-F-P053-018}{m_{n,1}=m_{n,5}=0} now make both integrals in \eqref{eq:5.20} zero. The
axial-viscosity term uses the order-\NSi{NS-F-P053-019}{n-1} axial-flux integral, which is also zero, including when that
order is zero.
\end{EESegment}

\begin{EESegment}{EE-TGT-P053-005}
Integrating the conservative equations now gives, one order at a time,
\NSFormula{NS-F-P053-020}{display}{5.21}
\begin{equation}\label{eq:5.21}
\begin{aligned}
\int_0^\infty r^2\mathscr R_{\theta,n}\,dr
&=\partial_t\int_0^\infty r^2u_{\theta,n}\,dr
 +\partial_z\int_0^\infty r^2\sum_{i+j=n}u_{z,i}u_{\theta,j}\,dr
 -\partial_{zz}\int_0^\infty r^2u_{\theta,n-1}\,dr
 =0,\\
\int_0^\infty r\mathscr R_{z,n}\,dr
&=\partial_t\int_0^\infty ru_{z,n}\,dr
 +\partial_z\int_0^\infty r\left(\sum_{i+j=n}u_{z,i}u_{z,j}+p_n\right)\,dr
 -\partial_{zz}\int_0^\infty ru_{z,n-1}\,dr
 =0.
\end{aligned}
\tag{5.21}
\end{equation}
Here \NSi{NS-F-P053-021}{\mathscr R_{j,n}=q^{-A-1+\lambda_n}\mathfrak r_{j,n}}, with
\NSi{NS-F-P053-022}{j\in\{\theta,z\}}, is the order-\NSi{NS-F-P053-023}{n} tangential residual in physical variables.
When \NSi{NS-F-P053-024}{n=1}, the last integral in \eqref{eq:5.21} means the convergent integral of
\NSi{NS-F-P053-025}{u_{\theta,0}-P} from
\end{EESegment}

\NSPage{054}%
\begin{EESegment}{EE-TGT-P054-001}
\eqref{eq:5.19}, and the differentiation is applied to that convergent, subtracted integral. Dividing out the common
physical powers gives
\NSFormula{NS-F-P054-001}{display}{none}
\[
\int_0^\infty R^2\mathfrak r_{\theta,n}\,dR=0,
\qquad
\int_0^\infty R\mathfrak r_{z,n}\,dR=0.
\]
Equation~\eqref{eq:5.9} can therefore also be written as
\NSFormula{NS-F-P054-002}{display}{none}
\[
\begin{aligned}
\mathcal T_{n,\theta}(R)
&=-R^{-2}\int_0^R\varrho^2\mathfrak r_{\theta,n}(\varrho)\,d\varrho
 =R^{-2}\int_R^\infty\varrho^2\mathfrak r_{\theta,n}(\varrho)\,d\varrho,\\
\mathcal T_{n,z}(R)
&=-R^{-1}\int_0^R\varrho\mathfrak r_{z,n}(\varrho)\,d\varrho
 =R^{-1}\int_R^\infty\varrho\mathfrak r_{z,n}(\varrho)\,d\varrho.
\end{aligned}
\]
This is the representation for the higher-order stress itself, with its sign retained.
\end{EESegment}

\begin{EESegment}{EE-TGT-P054-002}
\textbf{Step 5. Find the stress support and bound the order-one term at the outer edge.}
When \NSi{NS-F-P054-003}{n\geq2}, the residual itself is zero after \NSi{NS-F-P054-004}{X_+}. Each same-order term
contains a positive-order factor, and axial viscosity acts on order \NSi{NS-F-P054-005}{n-1>0}. The formula that
integrates backward from infinity then proves \NSi{NS-F-P054-006}{\supp\mathcal T_n\subset[X_-,X_+]}.

When \NSi{NS-F-P054-007}{n=1}, the only source left outside is
\NSi{NS-F-P054-008}{-\partial_{zz}u_{\theta,0}} in the angular equation. Beyond \NSi{NS-F-P054-009}{X_b}, the leading
profile is the \NSi{NS-F-P054-010}{z}-independent heat field, so \NSi{NS-F-P054-011}{\mathcal T_1} is still supported
inside the closed active annulus.
\end{EESegment}

\begin{EESegment}{EE-TGT-P054-003}
To get the bound near the edge, restrict further to the terminal collar where the exterior heat-flow cutoff is one.
There \NSi{NS-F-P054-012}{u_{\theta,0}=K(r,t)f_o}, where
\NSi{NS-F-P054-013}{K(r,t)=c_\infty s^{-A}\mathcal H(2(1-t)/s)} is the physical heat field from \eqref{eq:4.29}.
The smooth terminal multiplier \NSi{NS-F-P054-014}{f_o} is constant on its first plateau and equals one whenever
\NSi{NS-F-P054-015}{X\geq X_b}. Its exact construction is \eqref{eq:A.12}, shifted in the logarithmic radial coordinate.
In the local coordinate \NSi{NS-F-P054-016}{y=\log X}, the difference from one has the form
\NSi{NS-F-P054-017}{1-f_o=e^{-4/\delta_b^2}} times a smooth factor, where
\NSi{NS-F-P054-018}{\delta_b=\log(X_b/X)}. Direct differentiation gives
\NSFormula{NS-F-P054-019}{display}{none}
\[
a_z(\eta)=-\frac{2\eta}{L},
\qquad
b_z(\eta)=\frac{-2D\eta a_z+d\partial_\eta a_z}{L},
\qquad
\partial_{zz}(Kf_o)=Kq^{-2D}\left(a_z^2\partial_{yy}f_o+b_z\partial_yf_o\right).
\]
\end{EESegment}

\begin{EESegment}{EE-TGT-P054-004}
Every fixed \NSi{NS-F-P054-020}{\eta}-derivative of \NSi{NS-F-P054-021}{a_z,b_z} is bounded on
\NSi{NS-F-P054-022}{[-1,1]}. The derivatives \NSi{NS-F-P054-023}{\partial_yf_o} and
\NSi{NS-F-P054-024}{\partial_{yy}f_o} are bounded, respectively, by
\NSi{NS-F-P054-025}{Ce^{-4/\delta_b^2}\delta_b^{-3}} and
\NSi{NS-F-P054-026}{Ce^{-4/\delta_b^2}\delta_b^{-6}}. Backward integration improves the bound by three powers of
\NSi{NS-F-P054-027}{\delta_b} by Lemma~\ref{lem:A.9}, used with \NSi{NS-F-P054-028}{c=4} and
\NSi{NS-F-P054-029}{j=3,6}. The radial weights and the change between \NSi{NS-F-P054-030}{R} and
\NSi{NS-F-P054-031}{\delta_b} stay bounded and smooth on this fixed collar. After the stress power
\NSi{NS-F-P054-032}{q^{-A-1/2+\lambda_1}} has been factored out, Lemma~\ref{lem:A.9} gives
\NSFormula{NS-F-P054-033}{display}{5.22}
\begin{equation}\label{eq:5.22}
|\mathcal T_1|\leq Ce^{-4/\delta_b^2}\delta_b^{-3},
\qquad
|\partial^I\mathcal T_1|\leq C_Ie^{-4/\delta_b^2}\delta_b^{-N_I}.
\tag{5.22}
\end{equation}
\end{EESegment}

\begin{EESegment}{EE-TGT-P054-005}
The second estimate in \eqref{eq:5.22} comes from differentiating the backward integral. Each derivative adds only a
finite inverse power of \NSi{NS-F-P054-034}{\delta_b}. Near this edge, the inner factor in
\NSi{NS-F-P054-035}{\zeta} stays between two positive constants. This proves \eqref{eq:5.13} when
\NSi{NS-F-P054-036}{n=1}. At the inner edge, the same estimate follows from
\NSi{NS-F-P054-037}{\mathcal T_1=0} on \NSi{NS-F-P054-038}{[0,X_-]}. When
\NSi{NS-F-P054-039}{n\geq2}, every fixed derivative of \NSi{NS-F-P054-040}{\mathcal T_n/\zeta} is bounded because all
of these stresses have the same compact support inside the annulus.
\end{EESegment}
\end{proof}

\begin{EESegment}{EE-TGT-P054-006}
\subsection{The coefficient induction and the residual left by a finite sum}\label{sec:5.3}
The last two lemmas build and extend one coefficient after all lower orders have been finished. Lemma~\ref{lem:5.1} solves
the inner equations \eqref{eq:5.2} through \eqref{eq:5.6}. Lemma~\ref{lem:5.2} extends that solution and gives its support
properties while keeping those equations unchanged on \NSi{NS-F-P054-041}{0\leq X\leq X_-}. This subsection applies
the two lemmas repeatedly and estimates the residual left by every finite truncation in \eqref{eq:5.25}, including its
physical derivatives. The gain in decay must rise with the truncation order, while the loss in the exponent caused by
physical differentiation must not depend on that order. Equation~\eqref{eq:5.26} gives the exact bound.
\end{EESegment}
